\clearpage\section*{EV1-T2 -- Common test}\label{EV1-T2}
\addcontentsline{toc}{section}{EV1-T2 -- Common test}
\textbf{Practice paper}. Chapters 1 to 4.\par Duration: 2 h. No calculator. Marked out of 20. 14 November 2026.
\subsection*{Part 1 -- Fluency (6 points)}
Twelve questions, 0.5 point each. Only one correct answer; no justification required. An incorrect or missing answer scores 0, with no loss of points.
\begin{enumerate}
\item Calculate $4+2\times3$.\par A. $18$\quad B. $9$\quad C. $13$\quad D. $10$
\item Calculate $\frac12+\frac13$.\par A. $\frac{5}{3}$\quad B. $\frac{2}{5}$\quad C. $\frac{5}{6}$\quad D. $\frac{1}{6}$
\item Calculate $\dfrac{2^4}{2^{2}}$.\par A. $2$\quad B. $4$\quad C. $1$\quad D. $9$
\item A value of 50 increases by 20 \%. Final value?\par A. $70$\quad B. $40$\quad C. $10$\quad D. $60$
\item A value of 80 is obtained after a 20 \% decrease. Initial value?\par A. $100$\quad B. $96$\quad C. $64$\quad D. $60$
\item A 25 \% increase then a 20 \% decrease: overall coefficient?\par A. $1$\quad B. $\frac{21}{20}$\quad C. $\frac{4}{5}$\quad D. $\frac{9}{10}$
\item The solution of $2x+4=0$ is\par A. $2$\quad B. $8$\quad C. $-2$\quad D. $-8$
\item The slope between $(0;4)$ and $(2;10)$ is\par A. $3$\quad B. $-3$\quad C. $6$\quad D. $2$
\item The mean of $4, 6, 8$ is\par A. $18$\quad B. $6$\quad C. $7$\quad D. $4$
\item A sector of $90^\circ$ represents what fraction of the total?\par A. $\frac{1}{3}$\quad B. $\frac{3}{4}$\quad C. $\frac{1}{4}$\quad D. $\frac{1}{2}$
\item A probability is $\dfrac{2}{5}$. Its complement is\par A. $\frac{5}{2}$\quad B. $\frac{2}{5}$\quad C. $\frac{1}{5}$\quad D. $\frac{3}{5}$
\item For $f(x)=x^2+4$, $f(-2)$ is\par A. $6$\quad B. $16$\quad C. $8$\quad D. $0$
\end{enumerate}
\clearpage\subsection*{Part 2 -- Independent exercises (14 points)}
Every step must be justified. Exact results are expected unless stated otherwise.
\Needspace{12\baselineskip}
\subsection*{Exercise 1 -- Two growth processes (5 points)}
We compare two models for the output of a workshop. On day 0, each forecasts 30 items. Model A forecasts 10 extra items each day; model B forecasts a daily doubling. The outputs on day $n$ are therefore $a_n=30+10n$ and $b_n=30\times2^n$, for $n\in\N$. A running total here means the sum of the outputs of the days indicated.
\begin{enumerate}
\item State the type of each sequence and calculate the outputs on day 2. \hfill(1 pt)
\item Calculate the cumulative outputs for days 0 to 3 in each model. \hfill(2 pts)
\item Check that model B exceeds model A on day 1, then explain why this gap stays positive on all following days by comparing the daily increments. \hfill(2 pts)
\end{enumerate}
\Needspace{12\baselineskip}
\subsection*{Exercise 2 -- Imposed roots (5 points)}
We look for a polynomial function $f$ of degree 2, defined on $\R$. Its curve crosses the x-axis at the x-coordinates $-1$ and 5, and it passes through the point $(0;10)$. This information must be enough to determine the function.
\begin{enumerate}
\item Give its factorised form. \hfill(1 pt)
\item Determine its sign then solve $f(x)\leq0$. \hfill(2 pts)
\item Give its vertex form and its maximum. \hfill(2 pts)
\end{enumerate}
\Needspace{12\baselineskip}
\subsection*{Exercise 3 -- A controlled approximation (4 points)}
The function $f$ is defined on $\R$ by $f(x)=x^2+1$. We wish to approximate $f(3.02)$ using the tangent to its curve at the point with x-coordinate 3. The calculations are done without a calculator.
\begin{enumerate}
\item For $h\ne0$, calculate and simplify the rate of change between 3 and $3+h$. \hfill(1 pt)
\item Deduce $f\prime(3)$, then the equation of the tangent at the point with x-coordinate 3. \hfill(1 pt)
\item Use the tangent to approximate $f(3.02)$. Then calculate the error ``exact value minus approximation''. \hfill(2 pts)
\end{enumerate}
